\documentclass[12pt,a4paper]{article}

\usepackage[margin=1in]{geometry}
\usepackage[hidelinks]{hyperref}
\usepackage{setspace}
\usepackage[ngerman]{babel}
\usepackage{changepage}
\usepackage{newtxtext}

\setlength{\parskip}{0.5\baselineskip}
\setlength{\parindent}{0pt}

\onehalfspacing
\pagestyle{empty}

\begin{document}

Guten Tag allerseits,\\
Ich hoffe es geht Ihnen allen gut. Heute habe ich eine große Aufgabe, zwar ich soll in etwa 15 Minuten über Friedrich Nietzsches Werk \textit{„Die Geburt der Tragödie“} sprechen. Der Text beinhaltet so viele Geschichte und ist in sowohl der Literaturtheorie als auch der Philosophie so wichtig, dass ich kaum alles genügend wiedergeben könnte. Das heißt, ich werde gar nicht über Thomas Manns \textit{„Der Tod in Venedig“} sprechen, sondern nur wie Nietzsche zur Thematik \textbf{Lust an Leid und Leidenschaft} seinen Beitrag liefert.\par

Ich denke wir wissen alle wer Friedrich Nietzsche ist, aber trotzdem werde ich einen kurzen Überblick geben. Er wurde am 15. Oktober im Jahr 1844 in  Preußen geboren. Im Jahr 1864 begann er an der Universität Bonn Theologie zu studieren, aber hat ein Jahr später aufgehört an Gott zu glauben und hat daher das Studium gewechselt. Er studierte dann in Leipzig die Klassische Philologie. Jetzt sind wir an dem Punkt angelangt, der uns eigentlich interessiert. Weiß jemand vielleicht was \textit{„Klassische Philologie“} sein soll, denn ich glaube nicht Stellenbosch Universität bietet diesen Kurs. Laut Wikipedia heißt es:
\begin{adjustwidth}{1.5cm}{1.5cm}
    \begin{spacing}{1}
        \begin{center}
            \small
                \textit{Die Philologie ist die zusammenfassende Bezeichnung für die Sprach und Literaturwissenschaft einer Sprache oder eines Sprachzweiges und entstand im 16. Jahrhundert aus griechisch philologia, philos und logos, wörtlich „Liebe zur Sprache.“ Man unterscheidet die Altphilologie oder Klassische Philologie, die sich mit Altgriechisch und Latein befasst und die Neuphilologie, die Beschäftigung mit den modernen Sprachen.}
        \end{center}
    \end{spacing}
\end{adjustwidth}
Ein Beispiel der Neuphilologie wäre die Germanistik. Also er war eigentlich Sprachwissenschaftler. Er veröffentlichte im Jahr 1872 sein erstes Buch, und das war sogar \textit{„Die Geburt der Tragödie aus dem Geiste der Musik.“} Es gab solche Streit über das Buch, dass seine Studente bei der Universität Basel die Studien gewechselt haben, um nur nicht mehr Nietzsche als Lehrer zu haben. In jenem Buch handelt es von einer ähnlichen Thematik wie Lila Montag erwähnt hat: wir haben hier zwei Götterpersönlichkeiten, nämlich Apollon und Dionysos. Nietzsche versucht die Tragödie als Gattung und der Prozess wodurch sie entstand zu analysieren und erklären. Er weist dabei auf diese zwei Götter, da sie für die Kunst sehr wichtig sind. Es heißt im Buch:
\begin{adjustwidth}{1.5cm}{1.5cm}
    \begin{spacing}{1}
        \begin{center}
            \small
                \textit{Diese Namen entlehnen wir von den Griechen, welche die tiefsinnigen Geiheimlehren ihrer Kunstanschauung zwar nicht in Begriffen, aber in den eindringlich deutlichen Gestalten ihrer Götterwelt dem Einsichtigen vernehmbar machen. An ihre beiden Kunstgottheiten, Apollo und Dionysus, knüpft sich unsere Erkenntnis, dass in der griechischen Welt ein ungeheurer Gegensatz, nach Ursprung und Zielen, zwischen der Kunst des Bildners, der apollinischen, und der unbildlichen Kunst der Musik, also der des Dionysos, besteht: beide so verschiedene Triebe gehen neben einander her, zumeist im offnen Zwiespalt und sich gegenseitig zu immer neuen kräftigeren Geburten reizend, um in ihnen den Kampf jenes Gegensatzes zu perpetuiren.}
        \end{center}
    \end{spacing}
\end{adjustwidth}
Nun, was soll das alles? Dionysos ist Gott des, unter anderen, Theater, und Apollon der Dichtkunst. Die altgriechische Tragödie begann immer mit einem Chor, der halb gesungen halb gesprochen hat, daher ist Dionysos mit der Musik verbunden eben wie er mit der Tragödie als Gattung verbunden ist. Denken Sie mal, kann man eigentlich die Musik sehen? Wer Musik lesen kann vielleicht, aber wenn man Musik hört ist das eigentlich ohne etwas zu sehen. Deswegen sind sowohl die Musik und Dionysos mit einem Mangel an Form verknupft. Apollon ist aber mit der Plastik, also diese kleine Statuen die die Griechen immer gemacht haben, verbunden. Und, ein Dichter versucht immer etwas großartiges kurz und präzise zusammenzufassen, nicht war? Also ist Apollon der Gegenseiter des Dionysos, insofern er die Formlichkeit verkörpert.\par

Genau das ist eben die Hauptsache des Buchs. Es gab schon zu dieser Zeit eine ausgebreitete Theoretik zu Apollon und Dionysos, wie zum Beispiel Johann Joachim Winckelmann. Er hatte bereits im 18. Jahrhundert geschrieben über eine Statue von Apollon und wie diese die Idealen der Griechen verkörpert. Nietzsche aber hat identifiziert, dass der Mensch in sich beide das Apollinische und das Dionysische beinhaltet, erinnern sie sich daran dass ich früher gelesen haben von offnen Zwiespalt. Wenn Apollon mit Form verbunden ist, ist er auch mit Grenzen verbunden. Weil, man nur etwas identifizieren kann, halt Form geben kann, indem man es begrenzt. Indem ich sage \textit{„Ich bin Quin Barnard“} sage ich auch, \textit{„Ich bin nicht Doktor Isabel dos Santos.“} Das nennt man principium individuationis, oder das Prinzip des Individuum. Natürlich daher ist Dionysos grenzlos, und zersprengt er alle Grenzen und ist die formlose pure Wesen an sich. Wenn er die Ewigkeit ist, ist Apollon die Schranken der Zeit. Es entsteht fast ein Kampf zwischen den zwei Götter, weil sie die Gegensätze zu einander sind, und der Kampf findet in unseren Seelen statt.\par

Allerdings heißt das nicht, man solle irgendein Trieb bevorzügen. Nietsche schreibt:
\begin{adjustwidth}{1.5cm}{1.5cm}
    \begin{spacing}{1}
        \begin{center}
            \small
                \textit{„Alles Vorhandene ist gerecht und ungerecht und in beidem gleich berechtigt.“ Wir haben den Zustand der Individuation als den Quell und Urgrund alles Leidens, als etwas an sich Verwerfliches zu betrachten}
        \end{center}
    \end{spacing}
\end{adjustwidth}
Doch können wir es nicht als etwas Verwerfliches betrachten, denn ohne es, wären wir nicht und könnten wir überhaupt nichts betrachten. Der Rausch, die Grenzlosigkeiten ist für Nietzsche das Natürliche, denken Sie an dem All bevor das \textit{„Big Bang.“} Doch lediglich durch Individuation entsteht die Möglichkeit für Kunst und Musik, auch wenn diese dionysisch sind. Das Dionysische braucht das Apollinische zu existieren, und umgekehrt auch. Vor allem wichtig ist die Gleichgewicht. Wäre ich zum Beispiel zu apollinisch, da würde ich mich zu stark begrenzen, würde ich vergessen ich sei eigentlich Teil der Natur und in eine Lücke meiner Identität verfallen, weil ich alles wichtige abgegrenzt habe. Andererseits, wenn ich zu dionysisch wäre, da würde ich mich ganz und gar vergessen, da ich die Grenze meiner Identität auslöscht und würde ich denken es gibt nicht sowas wie ich.\par

Die Tragödie stellt für Nietzsche das alles dar. Die Tragödie fängt mit einem Chor an, also ist von formloser Musik und dem Dionysischen die Rede. Sie ist überhaupt Dionysisch weil sie ein Drama ist und wir in dem Theater sind. Der Tragische Held kämpft aber für seine Identität, normalerweise umsonst. Indem er für sich kämpft trotz der Hoffnungslosigkeit seiner Situation, trotz des Tragischen, verkörpert er den Versuch auf Ordnung und das Apollinische. Die Tragödie stellt unmittelbar ein Leid dar, aber laut Epikur sei nicht jede Leid zu vermeiden, weil sie uns in gewissermaßen vertrauter Natur ist. Vor allem beim Kampf zwischen diesen zwei Mächte, Triebe, in uns ist das Leid nicht zu vermeiden, weil alles ist.

Dankeschön

\end{document}
